\documentclass[11pt,a4paper]{article}
\usepackage[T1]{fontenc}
\usepackage[utf8]{inputenc}
\usepackage[polish]{babel}
\usepackage{lmodern}
\usepackage[margin=2.3cm]{geometry}
\usepackage{amsmath,amssymb}
\usepackage{booktabs}
\usepackage{enumitem}
\usepackage{tikz}
\usetikzlibrary{arrows.meta}
\setlist{nosep,leftmargin=*}
\setlength{\parindent}{0pt}
\setlength{\parskip}{5pt}
\renewcommand{\arraystretch}{1.1}

\title{Algorytmy optymalizacji dyskretnej, laboratorium 2}
\author{Jakub Retmańczyk}
\date{8 listopada 2026}

\begin{document}
\maketitle

\section*{Zadanie 1}
\textbf{Zmienne decyzyjne.} $x_{sl}\ge 0$ to liczba galonów paliwa, które firma $s\in S$ dostarcza na lotnisko $l\in L$.
Firma $s$ może dostarczyć najwyżej $p_s$ galonów, lotnisko $l$ potrzebuje $d_l$ galonów, a galon od firmy $s$ na lotnisku $l$ kosztuje $c_{ls}$ dolarów.

\textbf{Ograniczenia.}
\begin{align*}
\sum_{l\in L} x_{sl} &\le p_s, && s\in S && \text{(podaż firmy)}\\
\sum_{s\in S} x_{sl} &= d_l, && l\in L && \text{(zapotrzebowanie lotniska)}\\
x_{sl} &\ge 0, && s\in S,\ l\in L
\end{align*}

\textbf{Funkcja celu.} Minimalizacja łącznego kosztu dostaw w dolarach:
\[ \min \sum_{s\in S}\sum_{l\in L} c_{ls}x_{sl}. \]

\textbf{Wyniki.}
\begin{center}
\begin{tabular}{@{}lrrrrr@{}}
\toprule
galony & L1 & L2 & L3 & L4 & razem / podaż \\ \midrule
Firma 1 & 0 & 165\,000 & 0 & 110\,000 & 275\,000 / 275\,000 \\
Firma 2 & 110\,000 & 55\,000 & 0 & 0 & 165\,000 / 550\,000 \\
Firma 3 & 0 & 0 & 330\,000 & 330\,000 & 660\,000 / 660\,000 \\
\bottomrule
\end{tabular}
\end{center}

\begin{enumerate}[label=(\alph*)]
\item Minimalny łączny koszt wynosi \textbf{8\,525\,000\,\$}.
\item Tak, paliwo dostarczają wszystkie trzy firmy.
\item Firmy 1 i 3 sprzedają całe swoje paliwo, a Firmie 2 zostaje 385\,000 galonów.
\end{enumerate}

\section*{Zadanie 2}
\textbf{Zmienne decyzyjne.} $x_p\ge 0$ to tygodniowa produkcja wyrobu $p\in P$ w kilogramach.
Kilogram wyrobu $p$ sprzedaję po $v_p$ dolarów, materiał na niego kosztuje $u_p$ dolarów, a obróbka na maszynie $m\in M$ trwa $t_{pm}$ minut. Godzina pracy maszyny $m$ kosztuje $k_m$ dolarów, maszyna jest dostępna przez $h_m$ godzin w tygodniu, a popyt na wyrób $p$ wynosi $D_p$ kilogramów. Zysk z kilograma wyrobu to
\[ z_p=v_p-u_p-\sum_{m\in M} \frac{t_{pm}}{60}\,k_m. \]

\textbf{Ograniczenia.}
\begin{align*}
\sum_{p\in P} t_{pm}x_p &\le 60\,h_m, && m\in M && \text{(czas pracy maszyny w minutach)}\\
0\le x_p &\le D_p, && p\in P && \text{(popyt)}
\end{align*}

\textbf{Funkcja celu.} Maksymalizacja tygodniowego zysku w dolarach:
\[ \max \sum_{p\in P} z_p x_p. \]

\textbf{Wyniki.}
\begin{center}
\begin{tabular}{@{}lrrrr@{}}
\toprule
 & $P_1$ & $P_2$ & $P_3$ & $P_4$ \\ \midrule
zysk [\$/kg] & 4,20 & 5,50 & 4,55 & 3,75 \\
produkcja [kg] & 125 & 100 & 150 & 500 \\
popyt [kg] & 400 & 100 & 150 & 500 \\
\bottomrule
\end{tabular}
\end{center}

Maksymalny zysk wynosi \textbf{3632,50\,\$} tygodniowo. Wyroby $P_2$, $P_3$ i $P_4$ produkuję w ilości równej popytowi, a $P_1$ tylko 125 z 400 kg.
Cały dostępny czas wykorzystuje tylko maszyna $M_2$ (3600 min). $M_1$ pracuje 3525 min, a $M_3$ 2100 min. $P_1$ wypada najsłabiej, bo zajmuje maszynę $M_2$ aż 10 minut na kilogram.

\section*{Zadanie 3}
\textbf{Zmienne decyzyjne.} Dla okresu $j=1,\dots,K$ zmienne $x_j$ i $y_j$ to liczba sztuk wyprodukowanych normalnie i ponadwymiarowo, a $s_j$ to liczba sztuk w magazynie na koniec okresu.
Sztuka kosztuje $c_j$ lub $o_j$ dolarów, produkcja ponadwymiarowa jest ograniczona przez $a_j$, a zapotrzebowanie wynosi $d_j$. Produkcja normalna jest ograniczona przez $B=100$, magazyn mieści $W=70$ sztuk, przechowanie sztuki przez jeden okres kosztuje $h=1500$ dolarów, a na początku w magazynie jest $s_0=15$ sztuk.

\textbf{Ograniczenia.}
\begin{align*}
s_{j-1}+x_j+y_j-d_j &= s_j, && j=1,\dots,K && \text{(bilans magazynu)}\\
0\le x_j &\le B, && j=1,\dots,K && \text{(produkcja normalna)}\\
0\le y_j &\le a_j, && j=1,\dots,K && \text{(produkcja ponadwymiarowa)}\\
0\le s_j &\le W, && j=1,\dots,K && \text{(pojemność magazynu)}
\end{align*}

\textbf{Funkcja celu.} Minimalizacja łącznego kosztu produkcji i magazynowania w dolarach:
\[ \min \sum_{j=1}^{K}\left(c_jx_j+o_jy_j+h\,s_j\right). \]

\textbf{Wyniki.}
\begin{center}
\begin{tabular}{@{}lrrrr@{}}
\toprule
okres & 1 & 2 & 3 & 4 \\ \midrule
produkcja normalna & 100 & 100 & 100 & 100 \\
produkcja ponadwymiarowa & 15 & 50 & 0 & 50 \\
stan magazynu & 0 & 70 & 45 & 0 \\
zapotrzebowanie & 130 & 80 & 125 & 195 \\
\bottomrule
\end{tabular}
\end{center}

\begin{enumerate}[label=(\alph*)]
\item Minimalny łączny koszt wynosi \textbf{3\,842\,500\,\$}.
\item Produkcja ponadwymiarowa jest potrzebna w okresach 1, 2 i 4.
\item Magazyn jest pełny tylko w okresie 2.
\end{enumerate}
W okresie 2 produkcja jest najtańsza, dlatego firma produkuje wtedy na zapas.

\section*{Zadanie 4}
\textbf{Zmienne decyzyjne.} $x_{ij}\in\{0,1\}$ jest równe 1, jeśli łuk $(i,j)\in A$ należy do ścieżki.
Przejazd łukiem kosztuje $c_{ij}$ i trwa $t_{ij}$, a $T$ to limit czasu. Przyjmuję $b_{i^\circ}=1$, $b_{j^\circ}=-1$ i $b_v=0$ dla pozostałych miast.

\textbf{Ograniczenia.}
\begin{align*}
\sum_{j:(v,j)\in A}x_{vj}-\sum_{i:(i,v)\in A}x_{iv} &= b_v, && v\in N && \text{(zachowanie przepływu)}\\
\sum_{(i,j)\in A} t_{ij}x_{ij} &\le T && && \text{(limit czasu)}\\
x_{ij} &\in\{0,1\}, && (i,j)\in A
\end{align*}

\textbf{Funkcja celu.} Minimalizacja kosztu ścieżki:
\[ \min \sum_{(i,j)\in A} c_{ij}x_{ij}. \]

\textbf{Wyniki.}

\textbf{(a)} Najtańsza ścieżka to $1\to2\to3\to5\to7\to9\to10$. Kosztuje \textbf{13}, a jej czas wynosi 15.

\textbf{(b)} Mój egzemplarz ma 12 miast, $i^\circ=1$, $j^\circ=12$ i $T=20$. Na rysunku przy każdym łuku jest koszt (na czerwono) i czas (na niebiesko). Pogrubiona jest najtańsza ścieżka mieszcząca się w limicie.

\begin{center}
\begin{tikzpicture}[x=1.15cm,y=1.05cm,>=Stealth,
  v/.style={circle,draw,fill=gray!20,minimum size=6.5mm,inner sep=0pt,font=\small},
  e/.style={->,thin,gray!70!black},
  opt/.style={->,very thick,black},
  lab/.style={font=\scriptsize,fill=white,inner sep=1pt}]
\newcommand{\ct}[2]{\textcolor{red}{#1},\,\textcolor{blue}{#2}}
\node[v] (1) at (0,0) {1};
\node[v] (2) at (1.5,3.2) {2};
\node[v] (6) at (3.2,1.6) {6};
\node[v] (3) at (3.2,-1.2) {3};
\node[v] (8) at (5.2,-3) {8};
\node[v] (4) at (6,-1.2) {4};
\node[v] (7) at (6.8,1.6) {7};
\node[v] (10) at (7.4,0.1) {10};
\node[v] (11) at (10.4,-0.25) {11};
\node[v] (5) at (9,-1.5) {5};
\node[v] (9) at (10.2,-3) {9};
\node[v] (12) at (12,0) {12};
\draw[e] (1) -- node[lab] {\ct{2}{10}} (2);
\draw[e] (2) to[out=15,in=105] node[lab,pos=0.55] {\ct{3}{15}} (12);
\draw[opt] (1) -- node[lab] {\ct{4}{5}} (3);
\draw[opt] (3) -- node[lab] {\ct{3}{4}} (4);
\draw[opt] (4) -- node[lab] {\ct{3}{4}} (5);
\draw[opt] (5) -- node[lab] {\ct{4}{5}} (12);
\draw[e] (1) -- node[lab] {\ct{5}{6}} (6);
\draw[e] (6) -- node[lab] {\ct{4}{5}} (7);
\draw[e] (7) -- node[lab,pos=0.5] {\ct{5}{8}} (12);
\draw[e] (3) -- node[lab] {\ct{4}{7}} (8);
\draw[e] (8) -- node[lab] {\ct{2}{6}} (9);
\draw[e] (9) -- node[lab] {\ct{2}{6}} (12);
\draw[e] (4) -- node[lab] {\ct{6}{3}} (10);
\draw[e] (10) -- node[lab,pos=0.6] {\ct{5}{2}} (11);
\draw[e] (11) -- node[lab] {\ct{6}{2}} (12);
\draw[e] (2) -- node[lab] {\ct{1}{3}} (6);
\draw[e] (6) -- node[lab] {\ct{2}{2}} (3);
\draw[e] (5) -- node[lab] {\ct{1}{2}} (9);
\draw[e] (7) -- node[lab,pos=0.8] {\ct{2}{3}} (5);
\draw[e] (10) -- node[lab,left,pos=0.45] {\ct{2}{2}} (7);
\end{tikzpicture}
\end{center}

Bez limitu czasu najtańsza jest ścieżka $1\to2\to12$ o koszcie 5, ale jej czas to 25. Z limitem najtańsza jest $1\to3\to4\to5\to12$ o koszcie \textbf{14} i czasie 18. Tyle samo kosztuje $1\to6\to7\to12$ o czasie 19.

\textbf{(c)} Tak. Kontrprzykład to dwie ścieżki z 1 do 4 i $T=10$. Ścieżka $1\to2\to4$ kosztuje 2 i trwa 16, a $1\to3\to4$ kosztuje 10 i trwa 4. Relaksacja daje $x_{ij}=0{,}5$ na wszystkich łukach i koszt 6, a optimum całkowite wynosi 10. Przyczyną jest ograniczenie czasu, przez które macierz nie jest całkowicie unimodularna.

\textbf{(d)} Nie zawsze. Bez limitu czasu macierz jest całkowicie unimodularna, więc sympleks zwraca ścieżkę, ale może ona przekraczać limit. W egzemplarzu (a) jest to $1\to3\to10$ o czasie $20>15$. Przy kilku ścieżkach o tym samym koszcie metoda punktu wewnętrznego może też zwrócić rozwiązanie ułamkowe (\texttt{d\_remis.dat}).

\section*{Zadanie 5}
\textbf{Zmienne decyzyjne.} $x_{dz}$ to liczba radiowozów przydzielonych do dzielnicy $d\in D$ na zmianie $z\in Z$ (liczba całkowita).
Tabele z listy na ćwiczenia dają dolne i górne limity $l_{dz}$ i $u_{dz}$, a $r_z$ i $q_d$ to minimalne liczby radiowozów na zmianie i w dzielnicy.

\textbf{Ograniczenia.}
\begin{align*}
\sum_{d\in D} x_{dz} &\ge r_z, && z\in Z && \text{(minimum dla zmiany)}\\
\sum_{z\in Z} x_{dz} &\ge q_d, && d\in D && \text{(minimum dla dzielnicy)}\\
l_{dz}\le x_{dz} &\le u_{dz}, && d\in D,\ z\in Z && \text{(limity z tabel)}\\
x_{dz} &\in\mathbb{Z}, && d\in D,\ z\in Z
\end{align*}

\textbf{Funkcja celu.} Minimalizacja łącznej liczby radiowozów:
\[ \min \sum_{d\in D}\sum_{z\in Z} x_{dz}. \]

\textbf{Wyniki.}
\begin{center}
\begin{tabular}{@{}lrrrrr@{}}
\toprule
 & zmiana 1 & zmiana 2 & zmiana 3 & razem & minimum \\ \midrule
$p_1$ & 2 & 7 & 5 & 14 & 10 \\
$p_2$ & 3 & 6 & 7 & 16 & 14 \\
$p_3$ & 5 & 7 & 6 & 18 & 13 \\ \midrule
razem & 10 & 20 & 18 & \textbf{48} & \\
minimum & 10 & 20 & 18 & & \\
\bottomrule
\end{tabular}
\end{center}

Potrzeba \textbf{48 radiowozów}. Mniej być nie może, bo minima dla trzech zmian dają razem $10+20+18=48$.

\section*{Zadanie 6}
\textbf{Zmienne decyzyjne.} $y_{pq}\in\{0,1\}$ jest równe 1, jeśli w kwadracie $(p,q)\in F$ stoi kamera. $F$ to zbiór wolnych kwadratów, bo na kontenerze kamery postawić nie można.
$C$ to zbiór kwadratów z kontenerami, a $V(i,j)$ to wolne kwadraty, z których widać kontener $(i,j)$:
\[ V(i,j)=\{(p,q)\in F:\ (p=i\wedge|q-j|\le k)\ \vee\ (q=j\wedge|p-i|\le k)\}. \]

\textbf{Ograniczenia.}
\begin{align*}
\sum_{(p,q)\in V(i,j)} y_{pq} &\ge 1, && (i,j)\in C && \text{(każdy kontener jest widziany)}\\
y_{pq} &\in\{0,1\}, && (p,q)\in F
\end{align*}

\textbf{Funkcja celu.} Minimalizacja liczby kamer:
\[ \min \sum_{(p,q)\in F} y_{pq}. \]

\textbf{Wyniki.} Mój teren ma $8\times10$ kwadratów i 18 kontenerów. Rozwiązałem go dla $k=1,2,3$. Na rysunkach K oznacza kontener, a C kamerę.

{\small
\begin{center}
\begin{tabular}{ccc}
$k=1$, 17 kamer & $k=2$, 7 kamer & $k=3$, 6 kamer \\
\begin{minipage}[t]{0.3\textwidth}\centering
\begin{verbatim}
C K K C . C C C K .
C . . C . K K . . C
K . . K . . . C . K
. C C K . C . K . .
. K . . . K . . C .
. . . . C K . . K .
K C K . C . . . . C
. . . . K . . . . K
\end{verbatim}
\end{minipage} &
\begin{minipage}[t]{0.3\textwidth}\centering
\begin{verbatim}
. K K C . . C . K C
. . . . . K K . . .
K . . K . . . . . K
. . . K . C . K . .
C K . . . K . . . .
. . . . . K . . K C
K . K . C . . . . .
. . . . K . . . . K
\end{verbatim}
\end{minipage} &
\begin{minipage}[t]{0.3\textwidth}\centering
\begin{verbatim}
C K K . . . . . K .
. . . C . K K . . .
K . . K . . . . . K
. . . K . . . K C .
. K . . . K . . . .
. . . . . K . . K C
K C K . . . . . . .
. . . . K C . . . K
\end{verbatim}
\end{minipage}
\end{tabular}
\end{center}
}

Dla $k=1$ potrzeba 17 kamer, czyli prawie tyle, ile jest kontenerów, bo tylko jedna kamera widzi dwa kontenery naraz. Dla $k=2$ wystarczy 7 kamer, a dla $k=3$ o jedną mniej.

\end{document}
